% Background, application context, and rationale.
\subsection{Background and Rationale}
The survey phase of seabed mapping can be modeled as a coverage problem: a vehicle must plan a path so that the relevant area is observed by its sensors. Coverage path planning is a central problem in robotics and includes marine applications such as image mosaicing and inspection \cite{galceran2013coverage}. In a simplified survey model, coverage can be represented as cells observed by the vehicle, while more detailed bathymetric reconstruction remains outside the scope of the experiment.

Using multiple vehicles can improve survey efficiency because the mission area can be split across the fleet. However, marine multi-robot systems introduce coordination problems. Prior work identifies communication, task division, and task allocation as closely connected challenges when marine links are limited or intermittent \cite{martorell2024coordination}. Connectivity-aware multi-robot exploration likewise shows that communication quality can affect mission time and robot roles \cite{abuaisheh2022cara}.

For a large survey area or a time-limited mission, multiple vehicles can perform disjoint survey work in parallel; the benefit depends on whether the saved survey time outweighs the coordination overhead. A simple solution is \emph{pre-partitioning}: assign each vehicle a non-overlapping region before launch and let it finish that region without further task exchanges. This is a strong option when workload and vehicle progress are predictable and balanced, so we retain it as a communication-independent baseline rather than assume it is inadequate.

The particular case motivating online reallocation is an initially balanced assignment followed by an unexpected, temporary slowdown of one vehicle. The other vehicle may finish its own region while unfinished work remains in the slowed vehicle's region. A fixed partition cannot transfer that work, whereas an online policy may redistribute unstarted tasks if current task state can be exchanged. Intermittent connectivity makes that exchange uncertain; advance reservation might help, but can also add overhead or leave work committed to the slower vehicle. The controlled slowdown in this study represents workload imbalance, not a calibrated model of currents or vessel dynamics. The experiment therefore tests when predictive reservation is useful, including conditions where fixed pre-partitioning or a non-predictive longer queue is sufficient.

Multi-robot coordination also appears in search and rescue, inspection, and exploration with restricted communication, not only marine surveys \cite{bramblett2023epistemic}. Tree coverage with a maintained communication constraint offers a contrasting terrestrial use case \cite{sinay2017tree}. These examples motivate distinguishing a coverage task from the networking mechanism used to coordinate it.

For a concrete marine application, Yang et al. motivate cooperative seabed-ROI surveys by offshore wind-farm site selection and construction \cite{yang2023cooperative}. Turnings, travel between survey sub-regions, and complete coverage matter in that application. Sousa et al. study aquatic sensing by a communicating USV swarm, illustrating a different marine networking objective: delivering measurements despite limited connectivity \cite{sousa2018aquatic}. Here the narrower practical difficulty is reallocating unfinished work when two surface vehicles progress unequally and cannot exchange current task state. The simulation will represent this task-level difficulty, not offshore conditions or sonar accuracy.
